\chapter*{Preface}
\addcontentsline{toc}{chapter}{Preface}
\markboth{Preface}{Preface}

Cuba is one of the few countries about which almost nobody writes carefully.
The literature divides into two libraries that do not read each other. One
records a small poor island that abolished illiteracy, built a health system
that outperformed its income by a wide margin, sent doctors to forty countries,
and did it under six decades of economic siege by the largest economy on earth.
The other records a one-party state that jailed its poets, criminalised its
markets, drove a fifth of its people into exile, and ended the period with a
grid that could not keep the lights on. Both libraries are accurate. Neither is
honest, because each is complete only by omission.

This book tries to hold both at once, and it does so for a practical reason
rather than a scholarly one. The last part of the book is a design: a proposed
socio-political and economic system for Cuba from 2026 forward. A design is
only as good as its account of the initial conditions. If you believe the first
library, you will design as though the only thing wrong is the embargo, and you
will produce a plan that fails the moment it meets a Cuban who wants to open a
shop. If you believe the second, you will design as though the only thing
wrong is socialism, and you will produce a plan that dismantles a public health
system it cannot replace, hands the assets to whoever is standing closest, and
reproduces in ten years the inequality that made 1959 possible. Both failures
have been run as experiments elsewhere. Neither needs a second trial.

So Parts I and II are not throat-clearing. They are the constraint set. Part I
covers 1959 to 1990, the years in which Cuba built its social state and paid
for it with a Soviet subsidy and a closed politics. Part II covers 1990 to
2026, the long adjustment to the loss of that subsidy: two full cycles of
liberalise-under-duress and recentralise-when-the-pressure-lifts, ending in the
worst material crisis the revolution has faced, worse in some dimensions than
the Special Period, and the largest emigration in Cuban history. Part III asks
what can be built on what is left.

What is left is more than the ruin makes it look. Cuba has a population that
can read, calculate and be trained quickly; a medical and biotechnological
establishment with real products and a reputation that is still worth money;
some of the best beaches, reefs and colonial cities in the hemisphere,
underexploited; eleven hundred kilowatt-hours of sunlight per square metre per
year falling on a country that currently burns imported fuel oil; a diaspora of
roughly two million people, many of them wealthy, most of them one flight from
home; and, unusually for a poor country, low violent crime and a functioning
public order. Those are the assets. The book's argument is that they compose:
that education plus connectivity plus a residence regime for mobile knowledge
workers is a second export track alongside tourism, that solar plus storage
converts the country's largest recurring hard-currency cost into a domestic
one, and that biotechnology is the one sector where Cuba already holds
something the world will pay for and has never been organised to sell.

The book's argument is also that none of it works without the boring part.
Every asset listed above has been available for twenty years and has not been
converted, and the reason is not the embargo alone. It is that Cuba lacks the
institutional machinery that turns an asset into an investment: enforceable
contracts, a currency that means something, a customs service that does not
take a cut, a court a foreigner can lose in fairly, a procurement system that
publishes what it bought. A country can have sunshine and doctors and beaches
and still be poor, and Cuba is the proof. So the longest chapters in Part III
are not the ones about solar panels. They are the ones about money,
about property, and about corruption, which is the failure mode every
transition of this type has actually died of.

One commitment in Part III is not up for argument in this book, and it should
be stated at the front rather than arrived at in the middle. The design keeps a
universal welfare state at Western European standard --- health, education,
pensions, disability, income support --- and everything else in Part III is
arranged around being able to afford it. This is not sentiment about the
revolution. It is the reading of the record in Parts I and II: the social state
is the one asset those sixty-seven years produced that is unambiguously worth
keeping, it is the thing that protects the Cubans most exposed if a transition
goes badly, and it is also an economic input rather than only a transfer, since
the education system is precisely what the second track and the biotechnology
sector are built out of.

What changes is not the benefit but who pays for it. Cuba's welfare state has
never been financed by Cubans. It was financed by Soviet terms of trade until
1990, by Venezuelan oil in the 2000s, and in the intervals by printing money
and by not maintaining anything, which is why it has come close to collapse
twice in living memory. An entitlement funded by an external rent lasts exactly
as long as the rent. The European systems this book takes as its standard are
financed by taxing a domestic base, which is why they survive changes of
government, of commodity price and of patron. So the proposal is a European
benefit structure on a European tax structure --- a broad-based, simply
administered tax system of a kind Cuba has never had --- and the cost of that
commitment, a tax burden well above the Latin American norm, is stated in the
chapters rather than hidden in them. A country that wants Nordic outcomes has
to collect Nordic revenue, and there is no version of this in which the
arithmetic is pleasant.

I write as a Cuban of the diaspora and a physicist by training, which shapes
the book in two ways worth declaring. It is why Part III is written as a set of
numbered institutional rules rather than as an essay: a design that cannot be
stated as rules is not a design, and rules can be argued with one at a time. It
is also why the book is not neutral about outcomes even where it tries to be
neutral about facts. It wants Cuba to be prosperous, open and its own. It does
not assume those three are automatically compatible, and where they conflict,
the text says so instead of choosing quietly.

Nothing here is a prediction. The prospect of the title is the older sense of
the word: a view of what lies ahead from where one is standing, drawn to scale.
